\documentclass[10pt,a4paper]{amsart}
\usepackage[margin=19mm]{geometry}
\usepackage{amsmath,amssymb,mathtools}
\usepackage[scr=rsfs,cal=euler]{mathalfa}
\usepackage{xfrac}
\usepackage[unicode,hidelinks,bookmarks=false]{hyperref}
\newcommand{\ie}{i.e.\ }
\newcommand{\cf}{cf.\ }
\newcommand{\resp}{respectively\ }
\newcommand{\Subsection}{\subsection}
\providecommand{\nobreakdash}{\nobreak\hbox{-}}
\setlength{\emergencystretch}{2em}
\multlinegap0pt
\usepackage{amssymb}
\usepackage{xcolor}
\usepackage{tikz}
\newtheorem{lemma}{Lemma}
\newcommand{\G}{\mathbb{G}}
\newcommand{\Z}{\mathbb{Z}}
\title{Stokes' Theorem on positively graded groups}
\author{Valentino Magnani, Francesca Tripaldi}
\date{Source: arXiv:2605.23594v1 (CC BY 4.0)}
\begin{document}
\maketitle

\noindent\emph{Source and licence.} Stokes' Theorem on positively graded groups. Version arXiv:2605.23594v1 (CC BY 4.0), distributed by arXiv under the Creative Commons Attribution 4.0 International licence, http://creativecommons.org/licenses/by/4.0/, as recorded in the arXiv licence metadata for this exact version and shown on the abstract page https://arxiv.org/abs/2605.23594v1. Full licence notice: \texttt{licenses/arxiv-2605.23594-CC-BY-4.0.txt}.\par\smallskip

\noindent\textit{Editorial scope.} Counted unit: the article's lemma less (source0.tex lines 1396--1398). The article's complete original proof of it is delivered with it (source0.tex lines 1399--1459, one \texttt{proof} environment of 61 lines). No mathematics is written, repaired or re-derived: the statement and the proof are the article's own lines, in the article's own notation, and no retained line is substituted or re-flowed.  No further source-local context is needed: every cross-reference of every retained line resolves inside the two delivered spans.  Nothing was omitted from any retained span, so the package carries no omission record.  No retained line contains a citation, so the wrapper carries no bibliography and no bibliography entry.  No retained line contains a figure environment, an image-inclusion command or drawing code, so no image is delivered and the package declares no asset dependency.  Editorial formatting only, in addition: an AMS \texttt{amsart} preamble that restates the article's own declarations and macros used by the retained lines, namely the environments \texttt{lemma} on separate counters, together with the standard packages those lines need.  The retained environments print in this document as Lemma 1 in order of appearance, whereas the article prints the same items with the numbering of its own full document.  The complete native source of the article as distributed in the arXiv e-print for this exact version is bundled unchanged as \texttt{sources/201224/source0.tex}.
\begin{lemma}\label{lem: deg boundary less}
Let $\Sigma\subset\G$ be a $C^1$ oriented manifold with boundary.  Then $\deg(\partial \Sigma)<\deg(\Sigma)$.
\end{lemma}

\begin{proof}
Let $k$ be the dimension of $\Sigma$. Let us consider $p\in\partial \Sigma$ such that $d_{\partial\Sigma}(p)=\deg(\partial \Sigma)$, let $(e_1,\ldots,e_k)$ be a basis of $T_p\Sigma$, such that $(e_1,\ldots,e_{k-1})$ is a basis of $T_p(\partial \Sigma)$. Whenever $1\le k'\le n$, we set
\[
I(k',n)=\{(j_1,j_2,\ldots,j_{k'})\in(\Z^+)^{k'}\mid  1\le j_1<j_1<\cdots<j_{k'}\le n \}
\]
and observe that there exists $\beta=(\beta_1,\ldots,\beta_{k-1})\in I(k-1,n)$ such that $\deg(\beta)=\deg(\partial \Sigma)$.
Consider now the matrix
\[
\left(
\begin{array}{cccc}
c_1^1 & c_2^1 & \cdots  & c_k^1 \\
\vdots & \vdots & \cdots & \vdots \\
c_1^n & c_2^n & \cdots  & c_k^n
\end{array}
\right)
=
C,
\]
where $e_i=\sum_{j=1}^n c_i^j X_j(p)$. We consider the $k$-tangent vector 
\[
\tau_{\Sigma}(p)=\sum_{\gamma\in I(k,n)}\tau_{\gamma}^{\Sigma}\; X_{\gamma}
=\sum_{\gamma\in I_{\beta}(k,n)}\tau_{\gamma}^{\Sigma} X_{\gamma}
+\sum_{\gamma\in I(k,n)\setminus I_{\beta}(k,n)}\tau_{\gamma}^{\Sigma} X_{\gamma}
\in\Lambda^k(T_p\Sigma)
\]
and we have defined
\[
I_{\beta}(k,n)
=
\left\{
(j_1,\ldots,j_k)\in (\mathbb N^+)^k
\;\middle|\;
1\leq j_1<\cdots<j_k\leq n,\;
\{\beta_1,\ldots,\beta_{k-1}\}
\subset
\{j_1,\ldots,j_k\}
\right\}.
\]
By contradiction, if we have 
\[
\tau_{\gamma}^{\bar\Sigma}=0
\qquad
\text{for all } \gamma\in I_{\beta}(k,n),
\]
then all rows of \(C\) are linear combinations of the rows
\[
(c_1^{\beta_1},\ldots,c_k^{\beta_1}), \; 
(c_1^{\beta_2},\ldots,c_k^{\beta_2}), \; 
\ldots,
(c_1^{\beta_{k-1}},\ldots,c_k^{\beta_{k-1}})
\]
of \(C\), therefore proving that the rank of $C$ is less than $k$. 
This conflicts with the fact that $\tau_\Sigma(p)\neq0$.
We have proved that there exists $\alpha\in I_{\beta}(k,n)$ such that 
$\tau^\Sigma_{\alpha}\neq0$, hence 
\[
d(\Sigma)\ge \deg(\alpha)>\deg(\beta)=\deg(\partial \Sigma),
\]
where the strict inequality holds since $\alpha$ contains all indexes of $\beta$.
This completes the proof.
\end{proof}

\end{document}
